\section{The Demonic Element in the Anti-Race}

\begin{quotex}
In this chapter from \emph{Sintesi di dottrina della razza}, \textbf{Julius Evola} takes up the question of the mismatch between one's spiritual and biological race. It is hopeless to understand what Evola is saying in normal terms. ``Race'', for him, is a spiritual quality that seeks to be manifested. However, there are other forces in play that may impede this process. In this regard, and also in terms of the differences of historical periods, he is closer to a Rudolf Steiner than to more conventional views. Here he admits the existence of demonic spirits.
\end{quotex}


In relation to that, and to complete the argument, it is time to consider the following. As a follow-up to the preceding objection we could call attention back to the fact that, in reality, the types at this point are not so differentiated, to be able to establish that aforesaid principle of fidelity to oneself; in the second place, the doctrine in question does not seem to give any explanation of the fact there exist types that are dilacerated and affected by serious conflicts, so much so that not all are ``their own type'' and not all feel themselves to be ``in their own home''.

On the basis of the general principle that everything that appears here is the analogical reflection of a reality existing above; in order to explain such cases one must think of everything that the aforesaid freedom of the individual without roots can do and likewise of the action of special historical conditions and social collectives; but above all, one should suppose, for such cases, some corresponding prenatal circumstances. Close to the central force that led to a given human manifestation, lesser and divergent forces can also have influences which, nevertheless, precisely because they are weaker, were, in a manner of speaking, overwhelmed and led to create expressions corresponding to elements of a unfavorable and discordant, biologico-historical, ``horizontal'' heredity.

The cases in which the ``race of the soul'' and the ``inner vocation'' do not correspond to the race of the body, such as those of every romantic laceration, in the final analysis, are to be explained in such a way from the metaphysical point of view. Even modern psychology knows, at this point, about so called ``split personalities''. And the more the minor forces diverge from the central direction, the more we will have, as the effect, men in whom the physical is not in accord with the soul, in whom the spirit contrasts with the body and the soul, whose vocation does not correspond to race or caste, the ``personality'' is in disagreement with tradition, and so on.

In similar cases, ``classical'' ethics, shaped by the ancient Nordic-Aryan norm of life, evident in a still more distinct way in its active and creative aspect, because it requires that the various divergent elements of these natures must obey a single, iron law on the basis of an inner decision, that cannot fail in the face of a negative example: and, as we will see, activist race theory must arouse precisely such a decision in the greatest number of individuals of a nation, as the base for all the rest. To exalt the romantic instead, the tragic, anxious soul, always in search of new ``truths'', is essentially something from a civilization sick and damaged in its race. Calmness, style, clarity, command, discipline, power and the Olympic spirit are instead the points of reference for every formation of character and life in the Nordic-Aryan sense.

But if even in the world of causes and metaphysical meanings we must presume the existence of natures and vocations manifesting a diverse degree of unitariness, we must also think that not all civilizations, and not all periods, offer the same possibilities of incarnation and expression to any of the forces that tend to an earthly form of existence. It is said that in every birth two different heredities interfere. The earthly and the historical gather, in a type of bond, certain biological, anthropological, and in part even psychological elements, a tradition, eventually also a caste, a given point in time and place in space, etc. Now, there are civilizations in which everything ``is in order'', so to say, in which life, at its maximum, develops around a great unity and organicity of all these elements of ``horizontal'' heredity. Other civilizations are instead characterized by individualism, anarchy, the destruction of every limit, and every difference deriving from race, blood, caste, tradition, nationality. From what was said about the law of elective affinity and analogous correspondences that act in birth, obviously the civilizations of the former type are those that, in order to offer them circumstances and possibilities of adequate expression, attract unitary natures and pure and decisive forces. The chaotic civilizations of the latter type, for the same reason, will become instead, so to speak, the ``exact place'', or the meeting place on earth, of every ``otherworldly hysteric''.

This curious expression is the least alarming that one can use to give the meaning of the thing. In fact, on the metaphysical plane, hysteria, the internal contradiction, can appear only as the quality of those beings, ``who say no to being'', to a greater or lesser measure. But such a quality is exactly that which Christian theology attributes to ``demonic'' forces --- now meant in the current sense of the term or to the ``creatures of chaos'', the will to incarnation of those, wherever situations arise that, for analogous reasons, evoke them, have thus a meaning as precise as it is alarming, that cannot be developed here. Typology, physiognomy, a type of transcendental psychology in the whole of a race theory examination of the first and second degree applied to the most typical figures of revolutionary and still outer, popular leaders of the frontlines of worldwide subversion, both political-social as well as cultural and spiritual, could lead, in this regard, to some impressive results.

We are not saying however that such chaotic civilizations accommodate these forces exclusively: unitary natures in themselves can also arise, which however will be particularly at a disadvantage and, in order to hold firm and keep faithful to a vocation, which in such cases often has the sense of a true mission, are condemned to dissipate a quantity of energy to face the conflicts between soul and body, race and character, inner dignity and rank, etc., that are typical of such civilizations and that the make them, in a normal way, the homeland of all other vocations. But in these cases it is necessary not to forget \textbf{Seneca}'s words, who correctly understood some unhappy circumstances in which a higher spirit can find himself involved in the discomforts and dangers to which he is exposed, or who is in a risky mission or in the line of combat: the most courageous and worthy are chosen for such tasks, while the cowardly and weak can be left to the ``comfortable life''.

It is not necessary, in any case, to emphasize the importance that the preceding considerations, although unusual for the common mentality of modern man, have for the racial idea and, in general, for the philosophy of civilization, once these exceptional cases are put aside. If a destiny of millennia has led the West to circumstances in which it would be difficult to still find something truly pure, intact, differentiated and traditional, to establish firm limits again, by any means, even with the harshest, is an endeavor whose beneficial effects perhaps at first sight may not be perceptible to everyone, but that nevertheless will not be found lacking in subsequent generations, through the effect of the secret way joining the visible with the invisible and the world with the ``superworld''.


\flrightit{Posted on 2014-03-19 by Aeneas}

\begin{center}* * *\end{center}

\begin{footnotesize}
\begin{sffamily}
\texttt{andrew on 2014-03-23 at 17:40 said:}

Although I do not understand all of what is being said, nevertheless it touches me. I would appreciate some practical guidance on what can be done by those who do not feel themselves to be ``in their own home''.

\end{sffamily}
\end{footnotesize}

